\label{sec:prong1}

\subsection{The Legal Standard}

Prong 1 requires that the minority group be ``sufficiently large and
geographically compact to constitute a majority in a single-member district.''
\textit{Gingles}, 478 U.S. at 50.  Courts have consistently held that Prong 1
is a geographic question, not a mathematical one: the relevant inquiry is
whether the minority community's geographic distribution makes it possible to
draw a compact district in which the minority constitutes a majority.  A
minority group that is large in absolute terms but geographically dispersed
across a state fails Prong 1.

The Supreme Court has declined to establish a precise compactness threshold
for Prong 1, relying instead on expert testimony and visual inspection.
This has produced inconsistency across circuits and made Prong 1 a battleground
for competing expert witnesses.

\subsection{VRASection's Geographic Alignment Score}

The VRASection component of \texttt{redist}~\citep{dellaLibera2026vra}
produces, as a byproduct of its optimization, a \emph{geographic alignment
score} for each potential minority-opportunity district.  The alignment score
$A_d$ for a district $d$ is defined as:

\begin{equation}
  A_d = \frac{\sum_{t \in d} \text{CVAP}^{\text{minority}}_t}
             {\sum_{t \in \text{COI}(d)} \text{CVAP}^{\text{minority}}_t}
\end{equation}

where the denominator sums minority CVAP over all tracts in the community of
interest (COI) identified by the VRASection algorithm as the geographic core
of the minority community, and the numerator sums minority CVAP over the
tracts actually assigned to district $d$.

An alignment score of 1.0 would indicate that district $d$ exactly captures
the minority community's core; scores above 0.6 indicate substantial alignment;
scores below 0.4 indicate that the district does not meaningfully correspond
to the minority community's geographic distribution.

\subsection{Operationalizing a Prong 1 Threshold}

\begin{proposition}
A minority group satisfies Prong 1 if and only if there exists a geographically
contiguous, population-balanced district configuration in which:
(a) the minority constitutes a majority of CVAP, and
(b) the district's alignment score $A_d \geq 0.5$ relative to the
    VRASection-identified community of interest.
\end{proposition}

The 0.5 threshold is derived from \textit{Gingles}~I (478 U.S.\ at 50--51):
``sufficiently large and geographically compact to constitute a majority in
a single-member district.''  We operationalise this as the minority comprising
$\geq 50\%$ of CVAP in a contiguous census-unit cluster.  The threshold
requires that the proposed district capture at least half the minority
community's geographic core, ensuring that the district reflects a genuine
geographic concentration rather than an artificial assembly of dispersed
minority tracts.

All majority-minority districts generated by VRASection for the 2020 census
have alignment scores above 0.55, and districts in states where Prong 1 has
been litigated and accepted by courts have alignment scores averaging 0.71.
The 0.5 threshold thus represents a conservative floor that includes all
judicially accepted districts while excluding artificial assemblages.

\subsection{Running the Prong 1 Analysis}

The Prong 1 analysis is produced by:
\begin{verbatim}
redist analyze --state AL --year 2020 --version remedy \
  --types demographic --vra-mode
\end{verbatim}
The output includes, for each minority-opportunity district: minority CVAP
percentage, alignment score, and a geographic visualization of the district
overlaid on the minority-population heatmap.  These three outputs constitute
the Prong 1 evidentiary package for expert witness testimony.
